\section{Results}
\label{sec:results}

\subsection{Corruption test}

Table~\ref{tab:corruption} gives the main result. Matching by value alone caught 1\% of the
corrupted numbers. The 25 tables hold thousands of values, so a corrupted number almost always
rounds to some unrelated cell and passes. Matching by surrounding words caught 20\% overall. Where the audit did choose a cell it
was usually the right one, and 88\% of corruptions among those numbers were caught, but it
flagged 44 correct numbers whose words pointed at the wrong cell.

\begin{table}[t]
\centering
\caption{Corrupted numbers caught on the DockPrepBench Results section (307 decimals, each changed
by three units in its last digit).}
\label{tab:corruption}
\begin{tabular}{lll}
\toprule
Method & Corruptions caught & Correct numbers flagged \\
\midrule
Value only & 1\% & 0 \\
Context (audit) & 20\% (88\% where anchored) & 44 \\
Model proposes, code verifies & 134 of 134 anchored & 0 \\
\quad (Claude Sonnet, blind mode) & & \\
\bottomrule
\end{tabular}
\end{table}

With Claude Sonnet proposing addresses in blind mode, 134 numbers were anchored and every one of
their corruptions was caught, with no correct number flagged. The remaining numbers were not
anchored, either because the model abstained or because its address did not reproduce the
original value, so the test says nothing about them. Coverage, not accuracy, is the limit of this
configuration.

\subsection{Local models}

Blind mode was not reliable with the local models. Qwen2.5-14B gave the right address for 43
numbers and a wrong one for 205; Qwen2.5-32B for 68 and 139. A wrong address cannot migrate a
number unless its value happens to round to the text, but in an audit each one is reported as a
possible error that the author has to look at.

Shortlist mode removes most of the difficulty. Table~\ref{tab:shortlist} compares the cells chosen
by the Qwen models with Claude's. Qwen2.5-32B linked fewer numbers than 14B but agreed with Claude
more often, and where the two Qwen models chose the same cell they agreed with Claude on 109 of
110 numbers.

\begin{table}[t]
\centering
\caption{Shortlist (migration) mode with local models, compared with Claude Sonnet's blind-mode
addresses.}
\label{tab:shortlist}
\begin{tabular}{lll}
\toprule
Model & Numbers linked & Same cell as Claude \\
\midrule
Qwen2.5-14B & 177 & 123 of 132 (93\%) \\
Qwen2.5-32B & 129 & 114 of 116 (98\%) \\
14B and 32B agree & 120 & 109 of 110 (99\%) \\
\bottomrule
\end{tabular}
\end{table}

\subsection{Migrating the manuscript}

The section was migrated in two passes. The first, with Claude's blind-mode addresses, replaced
124 numbers. The second offered the 245 numbers still typed by hand to Claude in shortlist mode;
it chose a cell for 211 and ``none'' for 34, and \texttt{migrate} replaced 188 where the formatted
value reproduced the text. After both passes 312 of the 369 decimal numbers in the section come
from data. The text extracted from the compiled PDF was identical before and after each pass,
which confirms that every replacement reproduces what was written.

The audit also found 28 numbers that no generated table produces. These are values computed by
hand or in one-off scripts, and they are the numbers most at risk after a rerun; the fix is to
add them to a table, not to the text.

\subsection{Model benchmark}

\result{Table of right / wrong / caught wrong / abstain / false anchor for the eight local models in
both modes, from \texttt{bench\_results.csv}, and two or three sentences on the ranking.}
